\clearpage
\section*{Public Blog: Stop Cooling the Empty Air}
\addcontentsline{toc}{section}{Public blog}

Would you leave your refrigerator door open to cool the kitchen in summer? The first breath of cold air feels wonderful, but the motor keeps working as warmth rushes back in. Outdoor cooling faces a similar challenge: a chilly outlet is useful only if its relief reaches the people waiting nearby.

\subsection*{1. Follow the comfort, not just the thermometer}
A fan can make you feel cooler without lowering air temperature. An evaporative cooler lowers temperature while adding moisture. Our model follows temperature, humidity, air movement and radiant heat to the visitor, then combines them using Standard Effective Temperature, or SET. One service unit means reducing one person's SET excess by one degree for one hour. Dividing that relief by electricity tells us how much comfort each kilowatt-hour buys.

\subsection*{2. A small queue teaches a big street a lesson}
We compared a $60$-square-meter shaded ticket queue at Honolulu Zoo with an $8{,}382$-square-meter occupied corridor beneath the Fremont canopy in Las Vegas \cite{honolulu-zoo-entrance,fremont-canopy}. The climates differ, and the queue concentrates more visitor-time into each square meter.

In Honolulu, three fans deliver about \textbf{143 service units using 3 kWh}: roughly \textbf{48 units per kWh}. In Fremont, one fan in the busiest zone achieves about \textbf{17 units per kWh}, yet reaches only \textbf{0.33\% of visitor-time} across the corridor. An excellent efficiency ratio can still leave almost everyone outside its reach.

Now enlarge the queue and its crowd eightfold. The same three fans keep their efficiency, but coverage falls from 100\% to 12.5\%. Twenty-four fans restore coverage. Scaling up means matching equipment to people and space.

\subsection*{3. Decide what good service means before buying}
Our example policy asks for at least 80\% coverage and $1.5^\circ$C average relief, with limits on occupied space and electricity. It selects 266 fans plus 16 evaporative units in Fremont, and one fan plus one evaporative unit in Honolulu. These are plans for the modeled conditions, not universal shopping lists.

The revealing detail is in the queue: three fans already deliver \textbf{nearly $1.5^\circ$C of relief}. A slightly more flexible target changes the least-cost choice to fans alone. Agree the service target first, then check that the equipment delivers it where people actually stand. During severe heat, follow local advisories and provide breaks and access to cooler places \cite{nws-heat}.

\vspace{5pt}
\begingroup\setlength{\fboxsep}{6pt}
\noindent\fcolorbox{black}{gray!5}{\begin{minipage}{\dimexpr\linewidth-2\fboxsep-2\fboxrule\relax}
\small\textbf{3 Quick Takeaways for Smart Public Cooling}
\begin{enumerate}[leftmargin=1.6em,itemsep=2pt,topsep=4pt,parsep=0pt]
\item \textbf{Cool people where they are.} Put useful airflow where visitors wait and gather.
\item \textbf{Check coverage as well as efficiency.} A high ratio does not tell you how many people benefit.
\item \textbf{Set the target, then compare costs.} Climate, natural wind and the required relief can change the best choice.
\end{enumerate}
\end{minipage}}\par\endgroup
